\emph{President, Montreal.AI and Quebec.AI, Montréal, Canada}

\textbf{Summary}\\
Machine intelligence becomes economically transformative only when it is
continuously converted into verified, allocable and reinvestable work
rather than isolated benchmark scores (refs. 1-4). I formalize
AGIJobManager as a protocol that mediates and constrains a distribution
over partially observed task environments and show that the resulting
economy admits a multiscale statistical-mechanical description. At micro
scale, each job is a protocol-augmented constrained POMDP whose
admissible trajectories are shaped by legality, settlement,
pause-adjusted time and scarce validation bandwidth. At meso scale, the
outer decision to enter a job is an entropy-regularized control problem
whose optimizer is a Gibbs-Boltzmann law over effective job energies,
while inner execution is an activated process whose successful
completion rate follows an Eyring-Kramers-like law over a coarse
free-energy barrier (refs. 5-13). At market scale, jobs, agents,
collateral, compute and validator bandwidth form a dilute
grand-canonical ensemble with chemical-potential-like shadow prices.
Coupling these laws to an exergy-limited accumulation model yields the
central result: compute accumulation alone saturates, whereas verified
machine labour can compound into persistent growth only when knowledge,
compute, energy throughput and automation co-accumulate (refs. 14-18).
Protocol design therefore appears as a free-energy engineering problem:
lower the activation barriers to validated work, enlarge the accessible
space of reproducible solution paths, and route the resulting surplus
into the physical capital that expands civilization's free-energy
frontier.

\section{Overview}\label{overview}

The public Prime architecture is unusually amenable to exact
formalization because it separates premium discovery from authoritative
settlement. Its public specification describes a settlement-first kernel
that preserves escrow, bonds, validator review, disputes, challenge
windows, pause controls and solvency accounting, together with a
procurement-first discovery layer for premium jobs. Two further semantic
commitments matter mathematically: deadlines advance in pause-adjusted
effective time, and ENS or public-page hooks are explicitly best-effort
side effects rather than authoritative state transitions. Those design
choices imply that the protocol is neither merely a payments rail nor a
conventional benchmark harness. It is a mechanism that makes a latent
task economy economically legible.

The governing statement of this paper is precise:

\[
\text{AGIJobManager mediates and constrains a distribution over partially observed task environments;}
\] \[
\text{an agent learns a hierarchical policy --- whether to enter a job, and how to act within it ---}
\] \[
\text{to maximize expected net economic utility subject to protocol and resource constraints.}
\]

What follows is a figureless, theorem-driven development of that
statement. The architecture is resolved across four coupled levels.
First, the micro level formalizes each job as a protocol-augmented
constrained POMDP. Second, the mesoscopic level derives Gibbsian entry
from entropy-regularized choice. Third, the kinetic level models
successful completion as activated barrier crossing in a partially
observed rugged landscape. Fourth, the macro level closes the system
with exergy-limited capital accumulation. The point is not ornament by
physics language. The point is that the relevant structures ---
partition functions, quasi-potentials, activated flux, shadow prices and
free-energy descent --- arise naturally once verified machine labour is
treated as a constrained stochastic process embedded in scarce
reservoirs.

\section{Formal set-up}\label{formal-set-up}

\textbf{Definition 1 (Protocol-mediated task ensemble).} Let job \(j\)
expose observable metadata \(z_j\) and induce a latent task environment

\[
\mathcal{T}_j = (\mathcal{S}_j,\mathcal{A}_j,\mathcal{O}_j,T_j,\Omega_j,u_j,\mathcal{C}_j,H_j),
\]

where \(\mathcal{S}_j\) is the latent state space, \(\mathcal{A}_j\) the
raw action space, \(\mathcal{O}_j\) the observation space, \(T_j\) the
transition kernel, \(\Omega_j\) the observation kernel, \(u_j\) the
economic utility functional, \(\mathcal{C}_j\) the resource constraints
and \(H_j\) the stopping rule. The protocol augments \(\mathcal{T}_j\)
with

\[
\Pi_j = (D_j,\Sigma_j,\mathcal{L}_j,t_{\mathrm{eff},j},E_j),
\]

where \(D_j\) is the procurement operator, \(\Sigma_j\) the
authoritative settlement map, \(\mathcal{L}_j(a\mid p)\) the legality
indicator induced by public protocol state \(p\), \(t_{\mathrm{eff},j}\)
the pause-adjusted effective clock and \(E_j\) the non-authoritative ENS
or public-page side channel. The resulting job environment is therefore
a \textbf{protocol-augmented constrained POMDP} with augmented state

\[
\tilde s_t = (x_t,h_t,m_t,p_t,r_t),
\]

where \(x_t\) is task workspace state, \(h_t\) hidden evaluator state,
\(m_t\) market state, \(p_t\) public protocol state and \(r_t\) local
resource state.

This decomposition isolates the protocol's real role. AGIJobManager does
not fully determine the task distribution: job creators, tools, hidden
tests and evaluators shape the latent task family. The protocol instead
mediates and constrains that family by declaring admissible actions,
authoritative settlement, timing semantics, procurement stages and
resource prices. That distinction is decisive. It is why the correct
formal object is not a single environment but a structured distribution
over environments coupled to a settlement mechanism.

\textbf{Definition 2 (Hierarchical admissible policy).} The agent's
control law is a pair

\[
\Pi = (\pi_{\mathrm{enter}},\pi_{\mathrm{act}}),
\]

where \(\pi_{\mathrm{enter}}\) acts at job-arrival times on job metadata
and current portfolio state, while \(\pi_{\mathrm{act}}\) acts inside an
accepted job over histories measurable with respect to pause-adjusted
effective time. Admissibility requires legality with respect to
\(\mathcal{L}_j\), satisfaction of hard protocol constraints and
feasibility with respect to local resource budgets.

The objective is

\[
\max_{\Pi}\;\mathbb{E}_{j\sim\mathcal{D}}
\left[
U(j,\Pi)\right]
\]

subject to protocol legality, collateral requirements,
validator-capacity constraints, compute budgets and time constraints.
Here \(\mathcal{D}\) is the task distribution induced jointly by job
creators, market conditions and protocol rules.

\section{Gibbsian job entry from entropy-regularized
control}\label{gibbsian-job-entry-from-entropy-regularized-control}

The outer decision to enter a job is not a pure argmax. Real agents face
uncertainty, finite search, opportunity cost and the need to preserve
exploratory diversity. The correct problem is entropy-regularized
allocation over the candidate job set plus the outside option. Let
\(b_j\) denote the belief state induced by job metadata and agent
history, and let \(V_j(b_j)\) denote the continuation value of solving
job \(j\) under the optimal inner policy. Consider

\[
\max_{q \in \Delta(\mathcal{J}_t \cup \{\varnothing\})}
\left\{
\sum_{j \in \mathcal{J}_t} q_j \widetilde V_j
- \Theta \sum_{j \in \mathcal{J}_t \cup \{\varnothing\}} q_j \ln \frac{q_j}{\rho_j}
\right\},
\]

with prior measure \(\rho_j\), decision temperature \(\Theta>0\) and net
continuation value

\[
\widetilde V_j = V_j - \lambda_C c_j - \lambda_B b_j - \lambda_V v_j - \lambda_T \tau_j,
\]

where \(c_j\), \(b_j\), \(v_j\) and \(\tau_j\) denote compute cost,
collateral burden, validator-bandwidth burden and effective-time burden.

\textbf{Theorem 1 (Canonical entry law).} The optimizer is

\[
q_j^{\star} = \frac{\rho_j\exp[-\beta\varepsilon_j^{\mathrm{eff}}]}
{1 + \sum_{k \in \mathcal{J}_t}\rho_k\exp[-\beta\varepsilon_k^{\mathrm{eff}}]},
\qquad \beta = \Theta^{-1},
\]

with effective job energy

\[
\varepsilon_j^{\mathrm{eff}} = -V_j + \lambda_C c_j + \lambda_B b_j + \lambda_V v_j + \lambda_T \tau_j.
\]

The associated entry free energy is

\[
\mathcal{F}_{\mathrm{entry}}(b_t)
= -\Theta \ln\!\left[1 + \sum_{k\in\mathcal{J}_t}\rho_k\exp(-\beta\varepsilon_k^{\mathrm{eff}})\right].
\]

This is exact convex duality rather than analogy. Any protocol feature
that lowers effective job energy or enlarges the multiplicity of
feasible high-value trajectories increases participation and validated
flux. Any feature that traps collateral, injects settlement uncertainty
or compresses viable solution paths raises the effective energy
landscape and suppresses entry.

\section{Activated execution in a partially observed rugged
landscape}\label{activated-execution-in-a-partially-observed-rugged-landscape}

Once a job is accepted, execution is a barrier-crossing process on a
rugged, partially observed landscape. The relevant quantity is not
merely expected payout. It is the accessibility of a validating
trajectory. Let \(\xi\) be a coarse reaction coordinate through belief
space, artifact space and protocol state. Define an effective activation
free energy

\[
\Delta G_j^{\ddagger}(t) = \Delta H_j^{\ddagger}(t) - \Theta\,\Delta S_j^{\ddagger}(t) + \Lambda_j^{\ddagger}(t),
\]

where \(\Delta H_j^{\ddagger}\) is the irreversible burden needed to
reach a validating state, \(\Delta S_j^{\ddagger}\) is the logarithmic
multiplicity of successful trajectories consistent with current belief
and tool affordances, and \(\Lambda_j^{\ddagger}\) is a protocol-induced
barrier capturing legality restrictions, procurement stage, challenge
risk, verifier uncertainty and queueing on scarce evaluation resources.

\textbf{Theorem 2 (Activated completion law).} Under a coarse-grained
first-passage approximation, validated completion is activated, with
instantaneous completion hazard

\[
\kappa_j(t) = \kappa_j^0(t)\exp[-\beta\Delta G_j^{\ddagger}(t)],
\]

where \(\kappa_j^0\) is a pre-exponential factor representing local
reasoning speed, tool throughput and reversible search dynamics. The
deadline-completion probability is therefore

\[
P_j^{\mathrm{dl}} = 1 - \exp\left[-\int_0^{T_{\mathrm{eff},j}} \kappa_j(t)\,dt_{\mathrm{eff}}\right].
\]

The structure mirrors advanced chemical kinetics but remains
disciplined. The barrier is a coarse quasi-potential in utility units,
not a claim that every economic quantity is literally molar
thermodynamics. The distinction matters. Decision temperature \(\Theta\)
belongs to entropy-regularized control, whereas physical temperature
\(T\) belongs to literal thermodynamic free energies. The formal analogy
becomes literal only when the job itself concerns physical design or
experimental optimization in chemistry, materials, energy or
manufacturing.

Operationally, the theorem gives design levers. Better tools, reusable
workflows and transferable credentials increase path multiplicity and
reduce entropic bottlenecks. Clearer settlement, lower dispute ambiguity
and faster validator throughput reduce the institutional barrier
\(\Lambda_j^{\ddagger}\). Higher compute, better reasoning systems and
improved tool chains raise the pre-exponential factor. The protocol
therefore shapes not only reward but kinetics.

\section{Grand-canonical markets and reservoir
prices}\label{grand-canonical-markets-and-reservoir-prices}

Consider a dilute market observed over an interval in which interactions
among simultaneous entrants are dominated by shared reservoirs: compute,
collateral and validator capacity. Let \(n_j\) denote the occupancy of
job class \(j\), let \(g_j\) denote the degeneracy of accessible
microtrajectories, and let \(\mu_C\), \(\mu_B\), \(\mu_V\) and \(\mu_A\)
denote shadow prices for compute, collateral, validator capacity and
agent activity. The grand partition function is

\[
\Xi = \prod_j \sum_{n_j=0}^{\infty} \frac{1}{n_j!}
\left[g_j\exp\!\left(-\beta\left(\varepsilon_j + \mu_C c_j + \mu_B b_j + \mu_V v_j - \mu_A\right)\right)\right]^{n_j},
\]

and therefore

\[
\langle n_j \rangle = g_j\exp\!\left[-\beta\left(\varepsilon_j + \mu_C c_j + \mu_B b_j + \mu_V v_j - \mu_A\right)\right].
\]

\textbf{Theorem 3 (Maxwell-Boltzmann market law).} In the dilute regime,
job occupancy obeys Maxwell-Boltzmann statistics over effective energies
shifted by reservoir prices, and the associated grand potential is

\[
\Omega = -\Theta \ln \Xi.
\]

Contestability and interoperability enter twice. They enlarge \(g_j\),
the accessible multiplicity of high-quality trajectories, and they lower
effective reservoir prices by reducing lock-in, duplicated verification
and switching cost. Closed systems can still create rents, but they do
so by trapping potential rather than minimizing free energy.

\section{Nonequilibrium extension and free-energy
dissipation}\label{nonequilibrium-extension-and-free-energy-dissipation}

The canonical law is the static envelope. Real markets are driven, noisy
and path-dependent. A minimal nonequilibrium closure is a birth-death
process over occupancy vectors \(n=(n_1,\dots,n_K)\) with entry and exit
rates satisfying local detailed balance,

\[
\frac{w_j^+(n)}{w_j^-(n+e_j)}
= \exp\!\left[-\beta\,\Delta \Phi_j(n)\right],
\]

where \(\Phi(n)\) is an effective mesoscopic potential and \(e_j\) is
the unit vector in coordinate \(j\). The stationary law is then

\[
\pi(n) \propto \exp\!\left[-\beta\Phi(n)\right],
\]

with interaction terms absorbed into \(\Phi\) whenever congestion,
strategic crowding or auction feedback matter.

\textbf{Proposition 4 (Relative free-energy Lyapunov functional).} Let
\(p_t(n)\) evolve under the master equation induced by the above rates.
Then the relative free energy

\[
\mathscr{F}[p_t\Vert \pi] = \Theta\,D_{\mathrm{KL}}(p_t\Vert \pi)
\]

is non-increasing in time, with equality only at stationarity.

This is the right thermodynamic statement for a protocol market. The
system need not be in equilibrium. What matters is that the protocol
defines the effective potential relative to which deviations dissipate.
Frictional mechanisms --- queueing, disputes, validator scarcity,
collateral lock-up, poor portability --- are then not vague
inefficiencies. They are sources of excess dissipation that keep the
realized market farther from its low-free-energy stationary measure.

\section{Exergy-limited accumulation and the compute-only
ceiling}\label{exergy-limited-accumulation-and-the-compute-only-ceiling}

Let \(X_t\) denote exergy throughput or exergy-capture capacity, \(C_t\)
compute capital, \(R_t\) robotics and automated actuation capital, and
\(Q_t\) validated knowledge capital. A minimal macro closure is

\[
Y_t = A_t(\eta_t X_t)^{\alpha} C_t^{\beta} R_t^{\gamma} Q_t^{\delta},
\qquad \alpha,\beta,\gamma,\delta>0,
\]

with updates

\[
X_{t+1}=(1-\delta_X)X_t+s_XY_t, \qquad
C_{t+1}=(1-\delta_C)C_t+s_CY_t,
\]

\[
R_{t+1}=(1-\delta_R)R_t+s_RY_t, \qquad
Q_{t+1}=(1-\delta_Q)Q_t+s_QJ_t,
\]

where \(J_t\) is validated job flux and \(\eta_t\) is coordination
efficiency.

\textbf{Theorem 5 (Compute-only saturation).} If exergy throughput
\(X_t\), robotic capital \(R_t\) and validated knowledge \(Q_t\) are
bounded above while \(0<\beta<1\), then no policy with \(s_X=s_R=s_Q=0\)
can sustain asymptotically positive exponential growth of \(Y_t\).

\emph{Sketch.} Under the stated bounds, \(Y_t \le \bar K C_t^{\beta}\)
for some finite \(\bar K\). Then

\[
C_{t+1} \le (1-\delta_C)C_t + s_C\bar K C_t^{\beta},
\]

which is sublinear in \(C_t\). Hence \(\log C_t/t \to 0\), and therefore
\(\log Y_t/t \to 0\).

The positive result is the converse. When validated work lowers future
barriers through \(Q_t\) and reinvestment simultaneously deepens exergy
throughput, compute and robotics, the system can cross into a
persistent-growth phase. In that regime, machine intelligence ceases to
be a one-shot productive factor and becomes a catalytic
capital-formation engine.

\section{Thermodynamic ceiling and astrophysical
scaling}\label{thermodynamic-ceiling-and-astrophysical-scaling}

The phase structure of the model is more revealing than any one-period
welfare comparison. More intelligence alone is not sufficient. More
energy alone is not sufficient. Persistent expansion requires
institutional contestability, reliable settlement, barrier-lowering
knowledge accumulation and continued deepening of the physical base. If
protocol-mediated machine labour becomes sufficiently catalytic ---
lowering the barriers to building the very energy, compute and robotic
systems that extend its own reach --- then the economically relevant
ceiling is no longer the balance sheet of an isolated firm. It is the
accessible free-energy budget of industrial civilization.

That statement is thermodynamic rather than rhetorical. In the deep
limit, the upper bound is set by the exergy flux that automated industry
can capture, transform and route through its own expanding
infrastructure. A protocol that continuously converts machine
intelligence into verified work, verified work into knowledge, and
knowledge into enlarged productive capacity therefore does more than
allocate jobs. It creates a mechanism by which algorithmic cognition can
be compounded into energy systems, fabrication systems, logistics
systems and, in the asymptote, into the engineered capture of
increasingly large fractions of the surrounding astrophysical
free-energy field.

\section{Discussion}\label{discussion}

The central claim can now be restated with formal precision.
AGIJobManager mediates and constrains a distribution over partially
observed task environments; an agent learns a hierarchical policy ---
whether to enter a job and how to act within it --- to maximize expected
net economic utility subject to protocol and resource constraints. At
micro scale this is a protocol-augmented constrained POMDP. At meso
scale it is a Gibbsian allocation law over jobs, agents and scarce
reservoirs. At kinetic scale it is an activated barrier-crossing process
on a partially observed rugged landscape. At macro scale it is a
nonequilibrium growth process limited by exergy, capital formation and
barrier-lowering knowledge accumulation.

Several design implications follow.

First, authoritative settlement is foundational. If settlement can be
displaced by unauditable discretion, the state space ceases to be
mechanically legible and the free-energy description loses its anchor.
Second, procurement must allocate scarce high-value tasks by quality
discovery rather than mere reaction speed. Third, pause semantics matter
because time is part of the control geometry. Fourth, public identity or
publication layers should remain additive side channels unless they can
be guaranteed not to perturb settlement truth. Fifth, the objective must
be evaluated in net utility rather than nominal payout because
collateral, compute, energy and validator time are all genuine state
variables.

The broader interpretation is that open intelligence markets are
barrier-lowering machines. A well-designed protocol converts local acts
of successful reasoning into lower future activation barriers, wider
accessible state spaces and more reproducible conversion of exergy into
validated work. In that precise sense, protocol design becomes a branch
of applied nonequilibrium thermodynamics for machine civilization.

\section{Methods}\label{methods}

\subsection{Protocol semantics and state
decomposition}\label{protocol-semantics-and-state-decomposition}

The formalization is tied to public semantics rather than private
assumptions. The repository used here describes a two-layer architecture
with a settlement-first kernel and a procurement-first discovery layer,
pause-adjusted effective time, and ENS hooks that remain side effects
while settlement and disputes stay authoritative. Those features are
represented by \(\Sigma_j\), \(D_j\), \(t_{\mathrm{eff},j}\) and
\(E_j\).

The hidden state decomposition \(\tilde s_t=(x_t,h_t,m_t,p_t,r_t)\)
keeps protocol-governed quantities explicit. \(x_t\) is task state,
including artifacts and tool outputs; \(h_t\) is hidden evaluator state,
including hidden tests and preference structure; \(m_t\) is market
state, including rival entrants and queueing load; \(p_t\) is public
protocol state, including job status, balances and pause flags; \(r_t\)
is the agent's local resource state, including compute budget, liquidity
and time.

\subsection{Derivation of Theorem 1}\label{derivation-of-theorem-1}

Introducing a Lagrange multiplier for normalization, the entry
functional

\[
\mathcal{L}(q,\lambda) = \sum_j q_j\widetilde V_j - \Theta\sum_j q_j\ln\frac{q_j}{\rho_j} + \lambda\left(1-\sum_j q_j\right)
\]

has first-order condition

\[
\widetilde V_j - \Theta\left(\ln\frac{q_j}{\rho_j}+1\right) - \lambda = 0.
\]

Solving yields

\[
q_j = \rho_j\exp\!\left(\beta \widetilde V_j -1 - \lambda/\Theta\right),
\]

and normalization gives the canonical partition function. Writing
\(\varepsilon_j^{\mathrm{eff}}=-\widetilde V_j\) gives the form reported
in the main text.

\subsection{Activated completion law}\label{activated-completion-law}

The barrier decomposition

\[
\Delta G^{\ddagger} = \Delta H^{\ddagger} - \Theta\Delta S^{\ddagger} + \Lambda^{\ddagger}
\]

is a coarse-grained potential in utility units. It does not identify all
economic quantities with literal molar thermodynamic quantities. Rather,
it defines the quasi-potential controlling first-passage probability
under large-deviation asymptotics. When the underlying task is itself
physical --- for example, a materials, chemistry or energy optimization
problem --- the same notation can be directly coupled to literal
physical Gibbs free energies and chemical potentials.

\subsection{Nonequilibrium closure}\label{nonequilibrium-closure}

The master-equation extension assumes local detailed balance with
respect to an effective mesoscopic potential. More refined closures
could introduce exclusion effects, correlated entry, adversarial
strategy, auction-theoretic feedback or explicit many-body interactions
among agents. In those cases the exponential family survives but the
effective potential acquires interaction terms and the
entropy-production accounting becomes more elaborate.

\subsection{Data and code
availability}\label{data-and-code-availability}

No external empirical dataset was required for the theoretical
derivations. This manuscript is a formal contribution; any future
empirical calibration would need observed task arrivals, settlement
outcomes, validator queues, collateral usage, compute expenditure and
exergy costs.

\section{References}\label{references}

\begin{enumerate}
\def\labelenumi{\arabic{enumi}.}
\tightlist
\item
  MontrealAI. \emph{AGIJobManagerPrime --- next-generation sovereign AI
  labor protocol featuring autonomous agent discovery, game-theoretic
  job markets, and institutional-grade on-chain settlement}. GitHub
  repository (2026).
\item
  Puterman, M. L. \emph{Markov Decision Processes: Discrete Stochastic
  Dynamic Programming}. Wiley (1994).
\item
  Kaelbling, L. P., Littman, M. L. \& Cassandra, A. R. Planning and
  acting in partially observable stochastic domains. \emph{Artificial
  Intelligence} \textbf{101}, 99-134 (1998).
\item
  Sutton, R. S., Precup, D. \& Singh, S. Between MDPs and semi-MDPs: a
  framework for temporal abstraction in reinforcement learning.
  \emph{Artificial Intelligence} \textbf{112}, 181-211 (1999).
\item
  Shannon, C. E. A mathematical theory of communication. \emph{Bell
  System Technical Journal} \textbf{27}, 379-423 and 623-656 (1948).
\item
  Jaynes, E. T. Information theory and statistical mechanics.
  \emph{Physical Review} \textbf{106}, 620-630 (1957).
\item
  Altman, E. \emph{Constrained Markov Decision Processes}. Chapman \&
  Hall/CRC (1999).
\item
  Oliehoek, F. A. \& Amato, C. \emph{A Concise Introduction to
  Decentralized POMDPs}. Springer (2016).
\item
  Wynne-Jones, W. F. K. \& Eyring, H. The absolute rate of reactions in
  condensed phases. \emph{Journal of Chemical Physics} \textbf{3},
  492-502 (1935).
\item
  Kramers, H. A. Brownian motion in a field of force and the diffusion
  model of chemical reactions. \emph{Physica} \textbf{7}, 284-304
  (1940).
\item
  Landauer, R. Irreversibility and heat generation in the computing
  process. \emph{IBM Journal of Research and Development} \textbf{5},
  183-191 (1961).
\item
  Touchette, H. The large deviation approach to statistical mechanics.
  \emph{Physics Reports} \textbf{478}, 1-69 (2009).
\item
  Seifert, U. Stochastic thermodynamics, fluctuation theorems and
  molecular machines. \emph{Reports on Progress in Physics} \textbf{75},
  126001 (2012).
\item
  Parrondo, J. M. R., Horowitz, J. M. \& Sagawa, T. Thermodynamics of
  information. \emph{Nature Physics} \textbf{11}, 131-139 (2015).
\item
  International Union of Pure and Applied Chemistry. \emph{Gibbs
  energy}. IUPAC Gold Book (2025).
\item
  International Union of Pure and Applied Chemistry. \emph{Chemical
  potential}. IUPAC Gold Book (2025).
\item
  International Union of Pure and Applied Chemistry. \emph{Gibbs energy
  of activation}. IUPAC Gold Book (2025).
\item
  International Energy Agency. \emph{Energy and AI}. IEA (2025).
\end{enumerate}
